% 3.5 Historical echoes
\subsection{Historical echoes}
\label{sec:history}

None of this is new. The pattern of a missing default forcing a rethink shows
up repeatedly in history, sometimes in rich countries, and two cases help
separate what is specific to poverty from what is general.

\paragraph{Rubber without rubber trees.} In 1942 the United States found
itself without its most important strategic raw material. A committee chaired
by Bernard Baruch reported that the country had lost to Japan
\rubberLostShare{} of its pre-war source of natural rubber
\parencite{baruch1942rubber}, at a moment when every tyre, gasket and hose in
the war effort depended on it. The default, which was to buy rubber from
plantations in Asia, was gone. The reframing was to stop asking where else
rubber could be bought and to ask whether it could be made. Over the next few
years the government and industry built a synthetic rubber industry almost
from nothing, a programme its historians called ``a project that had to
succeed'' \parencite{herbert1985}.

\paragraph{Heat without oil.} Iceland offers a quieter, longer example. The
country has little coal or oil of its own and plenty of hot water underground.
Geothermal heating spread through the twentieth century, and in the 1960s
Reykjavík's municipal heating system was extended throughout the city
\parencite{ragnars1980}. Then the oil crises of 1973 and 1979 turned imported
heating oil from an inconvenience into a strategic weakness, and they changed
the country's energy policy \parencite{orkustofnun_heating}. Measured by the
volume of heated buildings, geothermal heat went from \icelandGeoStart{} in
1970 to \icelandGeoMid{} in 1980 and \icelandGeoEnd{} in 1985, while oil fell
from \icelandOilStart{} to \icelandOilMid{} and then \icelandOilEnd{}
\parencite{orkustofnun2020heating}. Today about \icelandGeoHeating{} of the
energy used to heat Icelandic homes comes from geothermal sources, and
district heating reaches \icelandDistrictHeating{} of the population
\parencite{orkustofnun_heating}. A sudden relative scarcity, expensive oil,
accelerated a shift that local resources already made possible.

\begin{counterpoint}
Both of these cases cut against the thesis as much as they support it. The
synthetic rubber programme was not frugal in any sense. It was scarcity met
with the full abundance of a wartime industrial state, its chemists, its oil
refineries and its capital. Iceland's shift relied on an unusual geological
endowment and on decades of public investment in drilling and pipes. What the
two cases share with the bottle light is the first two steps of the mechanism:
the default disappeared, and the question was reframed. What they do not share
is the poverty. They suggest that the removal of a default can push innovation
in rich places too, and that what turns that push into a solution is usually
capability, not deprivation.
\end{counterpoint}
